\documentclass[11pt,a4paper]{article}
\usepackage[margin=2.6cm]{geometry}
\usepackage{amsmath,amssymb,amsthm}
\usepackage{booktabs}
\usepackage[hidelinks]{hyperref}

\theoremstyle{definition}
\newtheorem{definition}{Definition}
\theoremstyle{plain}
\newtheorem{lemma}{Lemma}
\newtheorem{theorem}{Theorem}
\theoremstyle{remark}
\newtheorem{convention}{Convention}
\newtheorem{remark}{Remark}

\newcommand{\R}{\mathbb{R}}
\newcommand{\N}{\mathbb{N}}
\newcommand{\code}[1]{\texttt{#1}}

\title{Shared foundation for the formal verification of the\\ SAS-to-Python credit-risk migration}
\author{Blueprint \code{credit-risk-foundation}}
\date{10 November 2026}

\emergencystretch=3em
\begin{document}
\maketitle

\section{Purpose and scope}

Two independent teams translate the same pipeline to Lean: one from the original SAS code, one
from the Python port. A third then proves or disproves that the two translations behave the
same. This document is the only material the two translators share. It fixes \emph{types and
vocabulary}: what a number is, what a table is, and what a missing cell looks like.

It deliberately fixes \emph{no behaviour} of SAS or of Python. How a missing value compares,
what a division by zero returns, how a value is rounded to an integer, how ties are ranked:
each translator states these from its own source, in its own blueprint. If this document
decided them, both translations would inherit the same answer and the equivalence proof would
prove nothing.

All Lean declarations live in the namespace \code{Tabular}.

\section{Three strengths of result}\label{sec:tiers}

A translated program is written once, for an arbitrary number type $\alpha$ carrying the
operations of Definition~\ref{def:real_like}. The same Lean text is then read at three types.

\begin{center}
\begin{tabular}{@{}lp{3.6cm}p{9.4cm}@{}}
\toprule
Tier & Statement is about & What a proof establishes \\
\midrule
1. Generic & every $\alpha$ with the operations & The two programs apply the same operations
in the same order. It holds under any rounding, so the outputs agree bit for bit on every
input, provided each primitive is the same function in both runtimes (see
Remark~\ref{rem:primitives}). \\
2. Exact & $\alpha = \R$ & The two programs are mathematically equal, but may compute in a
different order. Each such result marks a place where floating-point outputs can differ. \\
3. Rounded & every $\alpha$ with the rounding laws of Definition~\ref{def:rounded_real} & A
quantitative bound: the outputs differ by at most a stated $\varepsilon$ in binary64. \\
\bottomrule
\end{tabular}
\end{center}

A fourth reading, $\alpha = $ Lean's \code{Float} (Definition~\ref{def:real_like_float}),
carries no theorem. It makes every translation executable, so that it can be run on real rows
and compared with what SAS or Python actually printed.

\begin{remark}\label{rem:primitives}
Addition, subtraction, multiplication, division and the square root are fixed by IEEE~754.
The exponential and the logarithm are not: two conforming runtimes may return results that
differ in the last bit. A Tier~1 theorem that passes through \code{exp} or \code{log} is
therefore bit-exact only up to that difference.
\end{remark}

\section{Numbers}

\begin{definition}[Real-like number type]\label{def:real_like}
A type $\alpha$ is \emph{real-like} when it carries: the constants $0$ and $1$; addition,
subtraction, multiplication, division and negation; the cast of a natural number into
$\alpha$; the reading of a decimal literal (a mantissa, a sign of the exponent and a decimal
exponent) as an element of $\alpha$; the relations $<$ and $\le$, each decidable; the four
unary functions $\exp$, $\log$, $\operatorname{sqrt}$ and $\operatorname{floor}$ from $\alpha$
to $\alpha$; and two Boolean tests, \code{isNaN} and \code{isFinite}. No law is required of any
of these operations. In Lean this is the class \code{Tabular.RealLike}.
\end{definition}

\begin{definition}[Absolute value]\label{def:abs}
For a real-like type $\alpha$ (Definition~\ref{def:real_like}) and $x \in \alpha$, the
absolute value of $x$ is $-x$ if $x < 0$ and $x$ otherwise. In Lean this is
\code{Tabular.RealLike.abs}.
\end{definition}

\begin{definition}[Boolean equality test]\label{def:eqb}
For a real-like type $\alpha$ (Definition~\ref{def:real_like}) and $x, y \in \alpha$, the
Boolean equality test of $x$ and $y$ is true exactly when $x \le y$ and $y \le x$. In Lean this
is \code{Tabular.RealLike.eqb}.
\end{definition}

\begin{definition}[The real numbers are real-like]\label{def:real_like_real}
The real numbers $\R$ are a real-like type (Definition~\ref{def:real_like}) with their usual
constants, field operations, casts and order; a decimal literal is read as the rational number
it denotes; $\exp$ is the real exponential; $\log$ is the natural logarithm of Mathlib
(\code{Real.log}); $\operatorname{sqrt}$ is the real square root of Mathlib (\code{Real.sqrt});
$\operatorname{floor}(x)$ is the greatest integer at most $x$, as a real number; \code{isNaN}
is constantly false and \code{isFinite} is constantly true. In Lean this is the instance
\code{Tabular.instRealLikeReal}.
\end{definition}

\begin{definition}[Machine floats are real-like]\label{def:real_like_float}
Lean's type \code{Float} of IEEE~754 binary64 numbers is a real-like type
(Definition~\ref{def:real_like}) with the machine operations: the constants $0$ and $1$,
\code{Float} addition, subtraction, multiplication, division and negation, \code{Float.ofNat}
as the cast, \code{Float.ofScientific} for decimal literals, the \code{Float} comparisons $<$
and $\le$, \code{Float.exp}, \code{Float.log}, \code{Float.sqrt}, \code{Float.floor},
\code{Float.isNaN} and \code{Float.isFinite}. In Lean this is the instance
\code{Tabular.instRealLikeFloat}. It is for evaluation only; no theorem is stated about it.
\end{definition}

\begin{lemma}[Operations of the real instance]\label{lem:real_ops}
For the real-like structure on $\R$ of Definition~\ref{def:real_like_real} and every real $x$:
$\exp x$ is \code{Real.exp x}, $\log x$ is \code{Real.log x}, $\operatorname{sqrt} x$ is
\code{Real.sqrt x}, $\operatorname{floor} x$ is the integer floor $\lfloor x \rfloor$ cast to
$\R$, the absolute value of Definition~\ref{def:abs} is $|x|$, \code{isNaN x} is false and
\code{isFinite x} is true.
\end{lemma}
\begin{proof}
Each of the first four equalities and the last two hold by definition of the instance. For the
absolute value, split on $x < 0$: then $|x| = -x$; otherwise $0 \le x$ and $|x| = x$.
\end{proof}

\begin{lemma}[Equality test on the reals]\label{lem:eqb_real}
For real numbers $x$ and $y$, with the real-like structure of
Definition~\ref{def:real_like_real}, the Boolean equality test of Definition~\ref{def:eqb} is
true if and only if $x = y$.
\end{lemma}
\begin{proof}
The test is true exactly when $x \le y$ and $y \le x$, which is $x = y$ by antisymmetry.
\end{proof}

\begin{convention}[Out-of-domain arguments]\label{conv:domain}
On $\R$, Mathlib's total functions return a conventional value outside their domain:
$x / 0 = 0$, $\log 0 = 0$, $\log x = \log |x|$ for $x < 0$, $\operatorname{sqrt} x = 0$ for
$x < 0$, and $\exp$ never overflows. None of these is what SAS or IEEE arithmetic does. A
translation must therefore never rely on them. Every division, logarithm, square root and
exponential in a translation is either (i) guarded in the Lean text by the same test the
source performs, with the source's own result on the failing branch, or (ii) covered by a
contract precondition stating that the argument is in the domain. Where the source tests
whether a result is finite (overflow), the translation applies \code{isFinite} to the result;
on $\R$ that test is always true, on \code{Float} it is the machine test.
\end{convention}

\section{Rounding laws}

Definition~\ref{def:rounded_real} restates, with the operations of
Definition~\ref{def:real_like}, the class \code{NumpyLean.Spec.RoundedReal} of the LogosNumpy
library (commit \code{043bf11}, file \code{proofs/NumpyLeanProofs/Rounded.lean}). That library
proves the laws for its bit-level model of IEEE binary64 arithmetic. The Logos server does not
load that library, so the binary64 instance is not re-checked here: a Tier~3 theorem is proved
for every type satisfying the laws, and its application to binary64 rests on the LogosNumpy
proof. No law is stated for $\exp$, $\log$, $\operatorname{floor}$ or decimal literals; a
Tier~3 contract that needs one states it as a hypothesis.

\begin{definition}[Real-like type with rounding laws]\label{def:rounded_real}
A \emph{rounded real-like type} is a real-like type $\alpha$ (Definition~\ref{def:real_like})
together with a value map $\operatorname{toReal} : \alpha \to \R$ and three real constants
$u$ (the unit roundoff), $\eta$ (the least positive subnormal) and $m$ (the least positive
normal), such that the following hold. Call $x \in \alpha$ \emph{finite} when
\code{isFinite x} is true.
\begin{enumerate}
\item $0 \le u \le 1/2$, $0 \le \eta$ and $\eta = 2\,u\,m$.
\item $0$ and $1$ are finite, $\operatorname{toReal} 0 = 0$ and $\operatorname{toReal} 1 = 1$.
\item If $x$ is finite then \code{isNaN x} is false, $-x$ is finite and
$\operatorname{toReal}(-x) = -\operatorname{toReal} x$.
\item Addition and subtraction: if $x$, $y$ and $x \pm y$ are finite, there is a real $\delta$
with $|\delta| \le u$ and
$\operatorname{toReal}(x \pm y) = (\operatorname{toReal} x \pm \operatorname{toReal} y)(1 + \delta)$.
\item If $x$, $y$ and $x + y$ are finite, then $\operatorname{toReal}(x + y) = 0$ if and only
if $\operatorname{toReal} x + \operatorname{toReal} y = 0$.
\item Sterbenz: if $x$ and $y$ are finite and
$\operatorname{toReal} y / 2 \le \operatorname{toReal} x \le 2 \operatorname{toReal} y$, then
$x - y$ is finite and
$\operatorname{toReal}(x - y) = \operatorname{toReal} x - \operatorname{toReal} y$.
\item Multiplication: if $x$, $y$ and $x y$ are finite, there are reals $\delta, \varepsilon$
with $|\delta| \le u$, $|\varepsilon| \le \eta / 2$, $\delta \varepsilon = 0$ and
$\operatorname{toReal}(x y) = \operatorname{toReal} x \cdot \operatorname{toReal} y \cdot (1 + \delta) + \varepsilon$.
If moreover $m \le |\operatorname{toReal} x \cdot \operatorname{toReal} y|$, there is a real
$\delta$ with $|\delta| \le u$ and
$\operatorname{toReal}(x y) = \operatorname{toReal} x \cdot \operatorname{toReal} y \cdot (1 + \delta)$.
\item Division: the same two statements as for multiplication, with $x / y$ in place of $x y$
and $\operatorname{toReal} x / \operatorname{toReal} y$ in place of the product, under the
additional hypothesis $\operatorname{toReal} y \ne 0$.
\item Casts: if the cast of $n \in \N$ is finite, there is a real $\delta$ with
$|\delta| \le u$ and $\operatorname{toReal}(n) = n (1 + \delta)$; if $n\,u \le 1$, the cast of
$n$ is finite and $\operatorname{toReal}(n) = n$.
\item Order: for finite $x$ and $y$, $x < y$ if and only if
$\operatorname{toReal} x < \operatorname{toReal} y$, and $x \le y$ if and only if
$\operatorname{toReal} x \le \operatorname{toReal} y$.
\item Square root: if $x$ is finite and $0 \le \operatorname{toReal} x$, then
$\operatorname{sqrt} x$ is finite and there are reals $\delta, \varepsilon$ with
$|\delta| \le u$, $|\varepsilon| \le \eta/2$, $\delta\varepsilon = 0$ and
$\operatorname{toReal}(\operatorname{sqrt} x) = \sqrt{\operatorname{toReal} x}\,(1 + \delta) + \varepsilon$.
\end{enumerate}
In Lean this is the class \code{Tabular.RoundedReal}.
\end{definition}

\begin{definition}[The real numbers satisfy the rounding laws]\label{def:rounded_real_real}
The real numbers, with the real-like structure of Definition~\ref{def:real_like_real}, are a
rounded real-like type (Definition~\ref{def:rounded_real}) with
$\operatorname{toReal}$ the identity and $u = \eta = m = 0$. In Lean this is the instance
\code{Tabular.instRoundedRealReal}.
\end{definition}

\begin{lemma}[The real instance is exact]\label{lem:rounded_real_exact}
For the rounded structure on $\R$ of Definition~\ref{def:rounded_real_real}:
$\operatorname{toReal} x = x$ for every real $x$, the unit roundoff $u$ is $0$ and the least
subnormal $\eta$ is $0$.
\end{lemma}
\begin{proof}
By definition of the instance.
\end{proof}

\begin{theorem}[Exact arithmetic when the roundoff is zero]\label{thm:exact_of_u_zero}
Let $\alpha$ be a rounded real-like type (Definition~\ref{def:rounded_real}) with $u = 0$ and
$\eta = 0$, and let $x, y \in \alpha$ be finite. If $x + y$ is finite then
$\operatorname{toReal}(x + y) = \operatorname{toReal} x + \operatorname{toReal} y$. If $x - y$
is finite then
$\operatorname{toReal}(x - y) = \operatorname{toReal} x - \operatorname{toReal} y$. If $x y$ is
finite then
$\operatorname{toReal}(x y) = \operatorname{toReal} x \cdot \operatorname{toReal} y$. If
$\operatorname{toReal} y \ne 0$ and $x / y$ is finite then
$\operatorname{toReal}(x / y) = \operatorname{toReal} x / \operatorname{toReal} y$.
\end{theorem}
\begin{proof}
Each law of Definition~\ref{def:rounded_real} gives $\delta$ with $|\delta| \le u = 0$, hence
$\delta = 0$, and for the product and the quotient $\varepsilon$ with
$|\varepsilon| \le \eta/2 = 0$, hence $\varepsilon = 0$. Substituting gives the four
equalities. So a Tier~3 bound specialises to a Tier~2 equality.
\end{proof}

\section{Tables}

A table here is the content of one CSV hand-off file of the pipeline, or of one dataset in
memory: named columns in a fixed order, and rows in a fixed order. This is the type at which
the two translations are compared. Inside a stage, each translator may use any intermediate
representation it finds faithful to its source.

\begin{definition}[Cell]\label{def:cell}
Over a number type $\alpha$, a \emph{cell} is either a numeric cell holding an optional element
of $\alpha$, or a text cell holding an optional string. The absent option stands for a missing
value (the SAS numeric missing value, a null in a Python data frame). In Lean this is the
inductive type \code{Tabular.Cell} with constructors \code{num} and \code{str}.
\end{definition}

\begin{definition}[Row]\label{def:row}
Over a number type $\alpha$, a \emph{row} is a list of pairs of a column name (a string) and a
cell (Definition~\ref{def:cell}), in column order. In Lean this is \code{Tabular.Row}.
\end{definition}

\begin{definition}[Table]\label{def:table}
Over a number type $\alpha$, a \emph{table} is a list of rows (Definition~\ref{def:row}), in
row order. In Lean this is \code{Tabular.Table}.
\end{definition}

\begin{definition}[Reading a column of a row]\label{def:row_get}
For a row $r$ (Definition~\ref{def:row}) and a column name $c$, the value of $r$ at $c$ is the
cell of the first pair of $r$ whose name is $c$, and is absent when $r$ has no such pair. In
Lean this is \code{Tabular.Row.get}.
\end{definition}

\begin{definition}[Writing a column of a row]\label{def:row_set}
For a row $r$ (Definition~\ref{def:row}), a column name $c$ and a cell $v$
(Definition~\ref{def:cell}): if some pair of $r$ has name $c$, the result is $r$ with the cell
of every pair named $c$ replaced by $v$; otherwise it is $r$ with the pair $(c, v)$ appended at
the end. In Lean this is \code{Tabular.Row.set}.
\end{definition}

\begin{lemma}[Reading after writing]\label{lem:row_get_set}
Let $r$ be a row (Definition~\ref{def:row}), $c$ a column name and $v$ a cell. Reading
(Definition~\ref{def:row_get}) column $c$ of the row obtained by writing
(Definition~\ref{def:row_set}) $v$ at $c$ gives $v$. For every column name $c' \ne c$, reading
$c'$ after that write gives the same result as reading $c'$ in $r$.
\end{lemma}
\begin{proof}
By induction on $r$. For the empty row the write appends $(c, v)$, whose name is $c$ and is not
$c'$. For a row $(d, w) :: r'$, split on whether $d = c$ and on whether a pair named $c$ occurs
in $r'$; in each case the first pair named $c$ after the write holds $v$, and pairs with other
names are unchanged and keep their positions relative to one another.
\end{proof}

\begin{convention}[What the table type does not decide]\label{conv:tables}
Numeric columns of both languages are cells over the same $\alpha$; a Python integer column is
represented by casts of integers. Printed formats (\code{BEST12.}, fixed decimals) are not
part of a table. Whether two tables are compared up to row order, and by which key, is stated
in each equivalence contract, not here.
\end{convention}

\end{document}
